% Reconstructed text from V5 PDF page 35; non-native source
\par\smallskip\noindent\rule{\textwidth}{0.4pt}\par{\sffamily\footnotesize Original V5 page 35 -- recovered text}\par\begin{flushleft}\footnotesize\ttfamily
\par\smallskip
\par\smallskip
\noindent Diagnostics and Serving Infrastructure for MAS 35\par
\par\smallskip
\par\smallskip
\par\smallskip
\noindent tool-call spans, reasoning traces, and intermediate outputs. These simulations model a de-\par
\noindent ployment where diagnostic buffers share GPU memory with KV cache blocks, representing\par
\noindent architectures where traces are embedded in the serving state rather than stored independently.\par
\par\smallskip
\noindent KV-cache eviction diagnostic impact. Table 7 reports the results for a 50-conversation, 100-\par
\noindent block cache with 3 agents (deterministic, seed=42, warm-up=5 conversations). Under LRU\par
\noindent eviction, 58.5\% of evictions removed blocks containing diagnostic content, while maintaining\par
\noindent 100\% cache hit rate (67/100 diagnostic blocks retained). FIFO showed similar behavior. A\par
\noindent diagnostic-aware policy achieved 58.2\% diagnostic eviction (a 0.3 pp reduction) but at substan-\par
\noindent tial serving cost: cache hit rate dropped to 88.4\%. The diagnostic-aware policy trades 11.6\par
\noindent percentage points of cache hit rate for complete diagnostic block retention — suggesting, in\par
\noindent this simulated setting, that naive policy augmentation is insufficient and supporting the need for\par
\noindent joint optimization (direction O1 in Section 7).\par
\par\smallskip
\noindent Table 7: KV-cache eviction simulation (illustrative; 50 conversations, 100-block cache, 3 agents,\par
\noindent seed=42, simulated CPU workload). Hit rate = hits/(hits+misses). Diag-bearing = diagnostic-classified\par
\noindent evictions / total evictions. A “diagnostic block” contains keyword-matched content (error messages,\par
\noindent stack traces, tool return values). Metrics are cumulative over the full run after warm-up (5 conversa-\par
\noindent tions).\par
\par\smallskip
\noindent  Policy Evictions Diag-bearing Hit rate Diag retained\par
\noindent  LRU 2,688 58.5\% 100\% 67/100\par
\noindent  FIFO 2,641 60.3\% 100\% 64/100\par
\noindent  Diag-Aware 2,658 58.2\% 88.4\% 100/100\par
\par\smallskip
\par\smallskip
\par\smallskip
\noindent Serving-system comparison. Table 8 reports a separate scenario: 20 conversations, 240 re-\par
\noindent quests, 5,000-token cache (deterministic, seed=42). Flat LRU and RadixTree achieve compara-\par
\noindent ble hit rates but dramatically different eviction counts, illustrating that infrastructure architec-\par
\noindent ture choice affects the frequency of potential diagnostic-evidence disruption events even when\par
\noindent aggregate serving performance is similar.\par
\par\smallskip
\noindent Table 8: Serving-system simulation (illustrative; 20 conversations, 240 requests, 5,000-token cache,\par
\noindent seed=42, simulated CPU workload). Hit rate = hits/(hits+misses). Utilization = occupied cache tokens\par
\noindent / total cache capacity. Metrics cumulative after warm-up.\par
\par\smallskip
\noindent  Policy Hit rate Utilization Mean TTFT Evictions\par
\noindent  LRU Flat 95.45\% 91.1\% 13.456 ms 208\par
\noindent  RadixTree 95.74\% 25.6\% 13.240 ms 0\par
\par\smallskip
\par\smallskip
\noindent  Three caveats apply. The two scenarios use different cache sizes and workloads and are\par
\noindent not directly comparable. The simulations use a simplified pipeline (simulated CPU workload,\par
\noindent not production GPU serving); the cache sizes are small relative to production deployments; and\par
\noindent the diagnostic-content classification uses keyword-based detection rather than actual attribution\par
\noindent value assessment.\par
\end{flushleft}
